% Part: first-order-logic
% Chapter: tableaux
% Section: identity

\documentclass[../../../include/open-logic-section]{subfiles}

\begin{document}

\olfileid{fol}{tab}{ide}

\olsection{\usetoken{P}{tableau} met \usetoken{S}{identity}: gelijkheid gebruiken}

Voor !!{tableau}s met !!{identity} hebben we extra afleidingsregels
nodig. Hieronder staan de regels voor $\eq$. De termen $t$, $t_1$ en
$t_2$ zijn gesloten: ze bevatten geen vrije variabelen.

\begin{defish}
\AxiomC{}
\RightLabel{$\eq$}
\UnaryInfC{\sFmla{\True}{\eq[t][t]}}
\DisplayProof
\hfill
\AxiomC{\sFmla{\True}{\eq[t_1][t_2]}}
\noLine
\UnaryInfC{\sFmla{\True}{!A(t_1)}}
\RightLabel{$\TRule{\True}{\eq}$}
\UnaryInfC{\sFmla{\True}{!A(t_2)}}
\DisplayProof
\hfill
\AxiomC{\sFmla{\True}{\eq[t_1][t_2]}}
\noLine
\UnaryInfC{\sFmla{\False}{!A(t_1)}}
\RightLabel{$\TRule{\False}{\eq}$}
\UnaryInfC{\sFmla{\False}{!A(t_2)}}
\DisplayProof
\end{defish}
Deze twee regels werken anders dan alle andere regels. Je kunt
$\TRule{\True}{\eq}$ en $\TRule{\False}{\eq}$ pas toepassen als er al
\emph{twee} !!{signed formula}s op de tak staan: zowel
$\sFmla{\True}{\eq[t_1][t_2]}$ als $\sFmla{S}{!A(t_1)}$.

\begin{ex}
Neem aan dat $s$ en $t$ gesloten termen zijn. Dan geldt
$\eq[s][t], !A(s) \Proves !A(t)$:
\begin{oltableau}
  [\sFmla{\False}{\formula{A}(t)}, just = \TAss
    [\sFmla{\True}{\eq[s][t]}, just = \TAss
      [\sFmla{\True}{\formula{A}(s)}, just = \TAss
        [\sFmla{\True}{\formula{A}(t)}, just={\TRule{\True}{\eq}[2, 3]}, close]
      ]
    ]
  ]
\end{oltableau}
Dit is het substitutieprincipe voor identiteit: identieke objecten mag
je in een formule voor elkaar vervangen. Het heet ook wel het principe
van Leibniz.

Met !!{tableau}s kunnen we ook bewijzen dat $\eq$ symmetrisch is. Dat
betekent hier dat $\eq[s_1][s_2] \Proves \eq[s_2][s_1]$:
\begin{oltableau}
  [\sFmla{\False}{\eq[s_2][s_1]}, just = \TAss
    [\sFmla{\True}{\eq[s_1][s_2]}, just = \TAss
      [\sFmla{\True}{\eq[s_1][s_1]}, just = {$\eq$}
        [\sFmla{\True}{\eq[s_2][s_1]}, just = {\TRule{\True}{\eq}[2, 3]}, close]
      ]
    ]
  ]
\end{oltableau}
Regel~$2$ is hier de eerste
!!{formula}~$\sFmla{\True}{\eq[s_1][s_2]}$ die
$\TRule{\True}{\eq}$ nodig heeft. Regel~$3$ is de tweede formule en
heeft de vorm $\sFmla{\True}{!A(s_2)}$. Lees
$!A(x)$ hier als $\eq[x][s_1]$. Dan is $!A(s_1)$ hetzelfde als
$\eq[s_1][s_1]$ en is $!A(s_2)$ hetzelfde als $\eq[s_2][s_1]$.

Met tableaus kunnen we ook bewijzen dat $\eq$ transitief is. Dat
betekent hier dat
$\eq[s_1][s_2], \eq[s_2][s_3] \Proves \eq[s_1][s_3]$:
\begin{oltableau}
  [\sFmla{\False}{\eq[s_1][s_3]}, just = \TAss
    [\sFmla{\True}{\eq[s_1][s_2]}, just = \TAss
      [\sFmla{\True}{\eq[s_2][s_3]}, just = \TAss
        [\sFmla{\True}{\eq[s_1][s_3]}, just = {\TRule{\True}{\eq}[3, 2]}, close]
      ]
    ]
  ]
\end{oltableau}
Ook hier heeft $\TRule{\True}{\eq}$ twee formules nodig. De eerste is
regel~$3$: $\sFmla{\True}{\eq[s_2][s_3]}$. Daarbij speelt $s_2$ de rol
van~$t_1$ en $s_3$ de rol van~$t_2$. De tweede formule heeft de vorm
$\sFmla{\True}{!A(s_2)}$ en staat op regel~$2$. Lees $!A(x)$ nu als
$\eq[s_1][x]$. Dan wordt $!A(s_2)$ de formule
$\eq[t_1][t_2]$ (dus regel~$2$). $!A(s_3)$ wordt de
!!{formula}~$\eq[s_1][s_3]$ in de conclusie.
\end{ex}

\begin{prob}
Maak voor elk van de volgende gevallen een gesloten !!{tableau}:
\begin{enumerate}
\item $\sFmla{\False}{\lforall[x][\lforall[y][((x = y \land !A(x))
      \lif !A(y))]]}$
\item $\sFmla{\False}{\lexists[x][(!A(x) \land
    \lforall[y][(!A(y) \lif y = x)])]}$,\\
  $\sFmla{\True}{\lexists[x][!A(x)] \land
    \lforall[y][\lforall[z][((!A(y) \land !A(z)) \lif y = z)]]}$
\end{enumerate}
\end{prob}

\end{document}
